Algorithmic redistricting and commission redistricting are not alternatives. They are complements. The algorithm provides what human commissioners cannot: a verifiable, deterministic baseline produced from public data without partisan input. The commission provides what the algorithm cannot: legitimate deliberation over VRA requirements, community of interest considerations, and political subdivision integrity---judgments that require contextual knowledge and public accountability.

Mode~(b) integration---algorithm-as-first-draft---achieves the most favorable balance between these contributions. It preserves commission authority over VRA and other legally required adjustments while making those adjustments transparent and verifiable. The Commission Adjustment Protocol specifies the process in sufficient detail to be implementable by any existing commission without statutory change; its adoption as a commission operating rule is a feasible near-term step for commissions in California, Colorado, Michigan, Arizona, and other states with independent redistricting authority.

The Iowa model demonstrates that non-partisan criteria-based redistricting can produce durable, cost-effective maps that command public legitimacy. Iowa's limitation---the residual human judgment in the LSA's map-drawing---is precisely what the \texttt{bisect} platform eliminates. An algorithmic Iowa-style process would retain Iowa's institutional strengths (nonpartisan staff, legislative approval-or-reject, criteria-based constraint) while replacing manual map-drawing with verifiable algorithmic output.

For redistricting commissioners and administrators reading this paper, the practical implication is direct: the \texttt{bisect} platform can be run today on 2020 census data to produce baseline maps for any state, demonstrating its functionality before the 2030 redistricting cycle. States that begin experimenting with Mode~(b) integration in 2026--2028 will be in the strongest position to implement it effectively in 2031. The investment in training, process design, and public communication can be spread over five years rather than compressed into the six months after census data release.

The political barrier to Mode~(b) adoption is modest: unlike federal legislation (which requires congressional majorities) or ballot initiatives (which require statewide campaigns), a commission operating rule requires only a commission vote. Commissioners appointed to produce neutral maps should find a Mode~(b) algorithmic baseline---which provides the most transparent possible demonstration that the commission's maps are not gerrymandered---congruent with their institutional purpose.
